% !TEX root = all_beamer.tex
% ==============================================================================
% HEADER LATEX PER PRESENTAZIONI BEAMER
% Corso di Machine Learning - Giorgio Gambosi
% Università di Roma "Tor Vergata"
% ==============================================================================

% ------------------------------------------------------------------------------
% CLASSE DOCUMENTO E METADATI
% ------------------------------------------------------------------------------
\documentclass[aspectratio=169, 10pt]{beamer}

\subtitle{Course of Machine Learning}
\author{Giorgio Gambosi}
\institute[Tor Vergata]{Master's Degree in Computer Science\\University of Rome ``Tor Vergata''}
\date{A.Y. 2025--2026}

% ------------------------------------------------------------------------------
% PACCHETTI MATEMATICI
% ------------------------------------------------------------------------------
\usepackage{amsmath, amssymb, amsfonts, mathtools, amsthm}
\usepackage{bm}              % Simboli matematici in grassetto
\usepackage{upgreek}         % Lettere greche dritte
\usepackage{derivative}      % Gestione semplificata delle derivate

% ------------------------------------------------------------------------------
% PACCHETTI GRAFICI E VISUALIZZAZIONE
% ------------------------------------------------------------------------------
\usepackage{graphicx}        % Inclusione immagini
\usepackage{xcolor}          % Gestione colori
\usepackage{xkcdcolors}      % Colori aggiuntivi stile xkcd
\usepackage{microtype}       % Miglioramenti tipografici

% ------------------------------------------------------------------------------
% PACCHETTI PER TABELLE
% ------------------------------------------------------------------------------
\usepackage{booktabs}        % Tabelle professionali
\usepackage{array}           % Estensioni per tabelle
\usepackage{multirow}        % Celle multi-riga
\usepackage{multicol}        % Layout multi-colonna

% ------------------------------------------------------------------------------
% TIKZ E PGFPLOTS
% ------------------------------------------------------------------------------
\usepackage{tikz}
\usepackage{tikz-network}
\usepackage[customcolors]{hf-tikz}
\usepackage{pgfplots}
\pgfplotsset{compat=1.18}
\usepgfplotslibrary{fillbetween}

% Librerie TikZ
\usetikzlibrary{arrows, arrows.meta, decorations.pathmorphing, decorations.pathreplacing}
\usetikzlibrary{decorations.markings, snakes, positioning, backgrounds, fit, shadows}
\usetikzlibrary{automata, patterns, calc, intersections, through, shapes}
\usetikzlibrary{shapes.geometric, shapes.arrows, scopes, chains, matrix}

% ------------------------------------------------------------------------------
% ALTRI PACCHETTI
% ------------------------------------------------------------------------------
\usepackage[skins]{tcolorbox}  % Box colorati
\usepackage{subcaption}         % Sottocaption per figure
\usepackage{listings}           % Inclusione codice sorgente

% ------------------------------------------------------------------------------
% TEMA E STILE BEAMER
% ------------------------------------------------------------------------------
\usetheme{Madrid}
\usecolortheme{whale}
\useinnertheme{rounded}
\usefonttheme{professionalfonts}
\setbeamertemplate{navigation symbols}{}  % Rimuove simboli di navigazione

% ------------------------------------------------------------------------------
% DEFINIZIONE COLORI PERSONALIZZATI
% ------------------------------------------------------------------------------
% Colori principali del tema
\definecolor{MainBlue}{RGB}{31,56,100}      % Blu principale (header, titoli)
\definecolor{AccentBlue}{RGB}{70,130,180}   % Blu accento (blocchi)
\definecolor{AlertRed}{RGB}{160,30,30}      % Rosso per alert
\definecolor{ExGreen}{RGB}{30,120,60}       % Verde per esempi
\definecolor{BlockBg}{RGB}{220,232,246}     % Sfondo blocchi
\definecolor{torgray}{RGB}{220,225,235}     % Grigio Tor Vergata

% Colori alternativi (Tor Vergata style)
\definecolor{TvDark}{RGB}{30,50,80}         % Blu scuro header
\definecolor{TvMid}{RGB}{70,100,140}        % Blu medio blocchi
\definecolor{TvLight}{RGB}{210,220,235}     % Blu chiaro sfondi
\definecolor{TvAlert}{RGB}{160,30,30}       % Rosso alert
\definecolor{TvGreen}{RGB}{20,100,50}       % Verde esempi

% Colori supplementari
\definecolor{torred}{RGB}{153,0,0}
\definecolor{torcyan}{RGB}{0,112,192}
\definecolor{cerise}{HTML}{DA3A62}
\definecolor{teal}{HTML}{029386}
\definecolor{slateblue}{HTML}{5B5EA6}
\definecolor{orange}{HTML}{F97306}

% Colori per nodi TikZ
\definecolor{nodeblue}{RGB}{101,140,196}    % xkcd Faded Blue

% ------------------------------------------------------------------------------
% CONFIGURAZIONE COLORI BEAMER
% ------------------------------------------------------------------------------
\setbeamercolor{frametitle}{bg=MainBlue, fg=white}
\setbeamercolor{title}{bg=MainBlue, fg=white}
\setbeamercolor{structure}{fg=MainBlue}

% Blocchi standard
\setbeamercolor{block title}{bg=AccentBlue, fg=white}
\setbeamercolor{block body}{bg=BlockBg}

% Blocchi alert
\setbeamercolor{block title alerted}{bg=AlertRed, fg=white}
\setbeamercolor{block body alerted}{bg=red!8}

% Blocchi esempio
\setbeamercolor{block title example}{bg=ExGreen, fg=white}
\setbeamercolor{block body example}{bg=green!7}

% ------------------------------------------------------------------------------
% CONFIGURAZIONE FONT BEAMER
% ------------------------------------------------------------------------------
\setbeamerfont{frametitle}{size=\normalsize,series=\bfseries}
\setbeamerfont{block title}{size=\small,series=\bfseries}

% ------------------------------------------------------------------------------
% STILI TIKZ PERSONALIZZATI
% ------------------------------------------------------------------------------
\tikzset{
  >=latex,  % Frecce stile LaTeX di default
  % Nodi per reti neurali
  node/.style={thick,circle,draw=myblue,minimum size=22,inner sep=0.5,outer sep=0.6},
  node in/.style={node,green!20!black,draw=mygreen!30!black,fill=mygreen!25},
  node hidden/.style={node,blue!20!black,draw=myblue!30!black,fill=myblue!20},
  node convol/.style={node,orange!20!black,draw=myorange!30!black,fill=myorange!20},
  node out/.style={node,red!20!black,draw=myred!30!black,fill=myred!20},
  % Connessioni
  connect/.style={thick,mydarkblue},
  connect arrow/.style={-{Latex[length=4,width=3.5]},thick,mydarkblue,shorten <=0.5,shorten >=1},
  % Stili numerati per facile mappatura
  node 1/.style={node in},
  node 2/.style={node hidden},
  node 3/.style={node out}
}

% ------------------------------------------------------------------------------
% AMBIENTI TEOREMI
% ------------------------------------------------------------------------------
\setbeamertemplate{theorems}[numbered]
\newtheorem{mytheorem}{Theorem}
\newtheorem{mydefinition}{Definition}

% ------------------------------------------------------------------------------
% CONFIGURAZIONE LISTINGS (CODICE)
% ------------------------------------------------------------------------------
\lstset{
  language=Python,
  basicstyle=\ttfamily\tiny,
  keywordstyle=\color{MainBlue}\bfseries,
  commentstyle=\color{gray},
  frame=single,
  breaklines=true,
}

% ==============================================================================
% MACRO MATEMATICHE PERSONALIZZATE
% ==============================================================================

% ------------------------------------------------------------------------------
% SPAZI CALLIGRAFICI (Calligraphic spaces)
% ------------------------------------------------------------------------------
\providecommand{\calX}{\mathcal{X}}         % Spazio input
\providecommand{\calY}{\mathcal{Y}}         % Spazio output
\providecommand{\calH}{\mathcal{H}}         % Spazio ipotesi
\providecommand{\calT}{\mathcal{T}}         % Training set
\providecommand{\calA}{\mathcal{A}}         % Algoritmo
\providecommand{\calR}{\mathcal{R}}         % Rischio
\providecommand{\calF}{\mathcal{F}}         % Famiglia di funzioni
\providecommand{\calL}{\mathcal{L}}         % Loss function

% Alias per compatibilità
\providecommand{\Rr}{\mathcal{R}}
\providecommand{\Tr}{\mathcal{T}}
\providecommand{\Hr}{\mathcal{H}}

% ------------------------------------------------------------------------------
% RISCHIO E PROBABILITÀ
% ------------------------------------------------------------------------------
\providecommand{\barR}{\overline{\mathcal{R}}}  % Rischio empirico
\providecommand{\EmpR}{\overline{\mathcal{R}}}  % Rischio empirico (alias)
\providecommand{\Rbar}{\overline{\mathcal{R}}}  % Rischio empirico (alias)
\providecommand{\pM}{p_M}                       % Probabilità modello
\providecommand{\pC}{p_C}                       % Probabilità corretta
\providecommand{\Prob}[2]{\mathbb{P}_{#1}\!\left[#2\right]}  % Probabilità

% ------------------------------------------------------------------------------
% FUNZIONI INDICATRICI
% ------------------------------------------------------------------------------
\providecommand{\indicator}[1]{\mathbf{1}\!\left[#1\right]}  % Funzione indicatrice
\providecommand{\ind}[1]{\mathbf{1}[#1]}                      % Versione breve

% ------------------------------------------------------------------------------
% OPERATORI MATEMATICI
% ------------------------------------------------------------------------------
\providecommand{\argmax}{\operatorname*{arg\,max}}  % Argomento del massimo
\providecommand{\argmin}{\operatorname*{arg\,min}}  % Argomento del minimo
\providecommand{\vcd}[1]{\operatorname{VCdim}(#1)} % VC dimension
\providecommand{\defeq}{\stackrel{\text{def}}{=}}  % Definito come

% ------------------------------------------------------------------------------
% VETTORI E MATRICI (Bold Math)
% ------------------------------------------------------------------------------
% Vettori minuscoli
\providecommand{\bx}{\mathbf{x}}            % Vettore x
\providecommand{\bz}{\mathbf{z}}            % Vettore z
\providecommand{\bw}{\mathbf{w}}            % Vettore peso
\providecommand{\btt}{\mathbf{t}}           % Vettore target
\providecommand{\wvec}{\mathbf{w}}          % Vettore peso (alias)
\providecommand{\bth}{\boldsymbol{\theta}}  % Parametro theta

% Matrici maiuscole
\providecommand{\bX}{\mathbf{X}}            % Matrice dati
\providecommand{\bZ}{\mathbf{Z}}            % Matrice Z
\providecommand{\Xmat}{\mathbf{X}}          % Matrice dati (alias)
\providecommand{\Hmat}{\mathbf{H}}          % Matrice H
\providecommand{\Smat}{\mathbf{S}}          % Matrice scatter generica
\providecommand{\bS}{\mathbf{S}}            % Matrice scatter
\providecommand{\Ii}{\mathbf{I}}            % Matrice identità
\providecommand{\bZero}{\mathbf{0}}         % Vettore zero
\providecommand{\bPsi}{\boldsymbol{\Psi}}   % Matrice Psi

% Medie (overline)
\providecommand{\ox}{\overline{\mathbf{x}}} % Media vettore x
\providecommand{\ow}{\overline{\mathbf{w}}} % Media vettore w
\providecommand{\oX}{\overline{\mathbf{X}}} % Media matrice X

% Matrici scatter specifiche (analisi discriminante)
\providecommand{\SW}{\mathbf{S}_W}          % Within-class scatter
\providecommand{\SB}{\mathbf{S}_B}          % Between-class scatter
\providecommand{\ST}{\mathbf{S}_T}          % Total scatter
\providecommand{\GX}{G_{\mathbf{X}}}        % Matrice Gram

% ------------------------------------------------------------------------------
% COMANDI DI FORMATTAZIONE TESTO
% ------------------------------------------------------------------------------
\providecommand{\term}[1]{\textbf{\color{accent}#1}}        % Termine importante
\providecommand{\halert}[1]{{\color{AlertRed}\textbf{#1}}}  % Alert rosso
\providecommand{\mlalert}[1]{{\color{AccentBlue}\textbf{#1}}}  % Alert blu
\providecommand{\mlemph}[1]{{\color{MainBlue}\textit{#1}}}  % Enfasi blu

% ==============================================================================
% FINE HEADER
% ==============================================================================

\newcommand{\blankpage}{
\clearpage
\thispagestyle{empty}
\null
\clearpage
}
